\section{The ARMOR Framework}

Building on our analysis, we introduce ARMOR (Anchor Rollout and Mixed Optimization for RL). 
This framework replaces the passive KL penalty with an active stabilization mechanism comprising two components, each explicitly targeting a failure mode identified in Section~\ref{sec:kl}.

\subsection{Anchor Rollout: Active Mode Retention}
To address the mode-seeking nature of reverse KL, where the model tends to collapse by forgetting diverse solutions present in the reference, we propose \textbf{Anchor Rollout}.
Instead of relying on loss penalties to implicitly retain modes, we \emph{actively} inject reference samples into the training batch.
Specifically, for each query $x$, we construct a hybrid response group $G$ during rollout:
\begin{equation}
\begin{gathered}
y_i \sim \pi_{\theta_{\text{old}}} (\cdot|x), \quad 1 \le i \le |G|-1,\\
y_\text{anc} \sim \pi_{\text{ref}}(\cdot | x), \quad \text{s.t. } R(x, y_\text{anc}) = 1.
\end{gathered}
\end{equation}
Here, the auxiliary response $y_\text{anc}$ serves as an ``anchor sample'' drawn from the reference policy $\pi_{\text{ref}}$ and is guaranteed to be correct\footnote{This positive-only anchor is critical, as including negative samples drove the model \textit{away} from the reference policy, which will lead to increased instability \citep{Deepseek-V3.2}.} ($R=1$) through rejection sampling (Algo.~\ref{alg:armor}, Line~\ref{line:reject}).
By including this anchor in every response group, we provide an explicit signal preventing the policy $\pi_\theta$ from drifting away from known correct solutions. 
As shown in Fig.~\ref{fig:teaser_perf}, \textbf{this simple intervention effectively stabilizes training} where standard KL fails.

From an implementation perspective, our continual training setting initializes $\pi_\mathrm{ref}$ from the best checkpoint of $\pi_\theta$, ensuring identical model architectures.
This design streamlines deployment:
As outlined in Algo.~\ref{alg:armor} (Line~\ref{line:rollout}), we can synchronize the reference parameters $\theta_\text{ref}$ to the inference engine alongside standard updates.
This enables hybrid sampling without the computational overhead of switching between distinct generation engines.

{
\renewcommand{\algorithmicthen}{}
\begin{algorithm}[t]
\begin{algorithmic}[1]
\STATE \textbf{Input:} Initial parameters $\theta^0, \theta_\text{ref} \leftarrow \theta^0$, Dataset $\mathcal{D}$, hyperparameters $B, \mu, \alpha, \tau$.

\FOR{step $t = 1, \dots, T$}
    \STATE \textcolor{gray}{// Phase 1: Anchor Rollout}
    \STATE Initialize batch $\mathcal{B} \leftarrow \emptyset$
    \WHILE{$|\mathcal{B}| < B$}
        \STATE Sample query batch $\mathcal X \sim \mathcal{D}$

        \STATE \textcolor{gray}{// Shared engine by switching parameter}
        \STATE Load $\theta^{t-1}$, sample on-policy groups $\mathcal{Y}_{\text{on}}$
        \STATE Load $\theta_\text{ref}$, sample off-policy groups $\mathcal{Y}_{\text{off}}$\label{line:rollout}
        
        \STATE Compute rewards $\mathcal R_{\text{on}}$ and $\mathcal R_{\text{off}}$

        \STATE \textcolor{gray}{// Rejection sampling}\label{line:reject}
        \FOR{$x \in \mathcal X$}
        \IF{$\text{Var}(\mathcal R_{\text{on}|x}) > 0, \max(\mathcal R_{\text{off}|x})=1$}
            \STATE Let $y_{\text{anc}} = \arg\max_{y\in \mathcal{Y}_{\text{off}|x}} R(x, y)$
            \STATE $G \leftarrow \mathcal{Y}_{\text{on}|x} \cup \{y_{\text{anc}}\}$
            \STATE $\mathcal{B} \leftarrow \mathcal{B} \cup \{(x,G)\}$ 
        \ENDIF
        \ENDFOR
    \ENDWHILE
    
    \STATE \textcolor{gray}{// Phase 2: Mixed Optimization}
    \FOR{iteration = $1,\dots,\mu$}
    \STATE Update $\theta^t$ by maximizing $\mathcal{J}(\theta)$, using the mixed IS ratio $r_{i,t}^\text{mix}$ (Eq.~\ref{eqn:mixed_is})
    \ENDFOR
    
    \STATE \textcolor{gray}{// (Optional) Reference Reset}\label{line:ref}
    \IF{$t \mod \tau = 0$}
        \STATE $\theta_\text{ref} \leftarrow \theta^t$
    \ENDIF
    % \STATE $\theta_\text{ref} \leftarrow \theta^t$ \textbf{if} $t \mod \tau = 0$ \quad\textcolor{gray}{// Reference Reset} 
\ENDFOR
\end{algorithmic}
\caption{ARMOR Training Framework}\label{alg:armor}
\end{algorithm}
}

\subsection{Mixed Optimization: Performance Ceiling}
While Anchor Rollout guarantees stability, our objective extends beyond collapse prevention to maximizing the model's asymptotic performance.
To achieve this, we first revisit the theoretical implications of our rollout strategy, observing that by injecting anchor samples, Anchor Rollout implicitly constructs \textbf{mixture policies} regularized by the reference:
\begin{equation}
\begin{gathered}
\pi^{\text{mix}}_\theta = \alpha\pi_\theta + (1-\alpha)\pi_{\text{ref}},\\
\pi^{\text{mix}}_{\theta_\text{old}} = \alpha\pi_{\theta_\text{old}} + (1-\alpha)\pi_{\text{ref}}
\end{gathered}
\end{equation}
where $\alpha \in (0, 1)$ is the mixing coefficient.
Since the data is generated from this mixture, the \textbf{principled optimization objective} should align with this structure.
Therefore, rather than optimizing $\pi_\theta$ in isolation, we reformulate the objective to optimize the mixture policy itself.

We implement this by substituting the standard probability terms in the policy gradient (e.g., DAPO) with their mixed variants, yielding the mixed Importance Sampling (IS) ratio:
\begin{equation}\label{eqn:mixed_is}
r_{i,t}^\text{mix}(\theta) = \frac{\alpha\pi_\theta(\cdot) + (1-\alpha)\pi_{\text{ref}}(\cdot)}{\alpha\pi_{\theta_{\text{old}}}(\cdot) + (1-\alpha)\pi_{\text{ref}}(\cdot)}.
\end{equation}
% \begin{equation}\label{eqn:mixed_is}
% r_{i,t}^\text{mix}(\theta) = \frac{\alpha\pi_\theta(y_{i,t}|x,y_{i,<t}) + (1-\alpha)\pi_{\text{ref}}(y_{i,t}|x,y_{i,<t})}{\alpha\pi_{\theta_{\text{old}}}(y_{i,t}|x,y_{i,<t}) + (1-\alpha)\pi_{\text{ref}}(y_{i,t}|x,y_{i,<t})}.
% \end{equation}
This formulation explicitly accounts for the reference policy's contribution to the data distribution.
Crucially, since $\pi_\text{ref}$ is fixed within each step, any performance gain achieved on the mixture target $\pi^{\text{mix}}_\theta$ strictly translates to improvements in the target policy $\pi_\theta$:
\begin{align*}
&\quad\mathbb E_{y\sim\pi_{\theta}^\text{mix}(\cdot|x), y'\sim\pi_{\theta_\text{old}}^\text{mix}(\cdot|x)}
[R(x,y) - R(x,y')] > 0 \\
&\Rightarrow \mathbb E_{y\sim\pi_{\theta}(\cdot|x), y'\sim\pi_{\theta_\text{old}(\cdot|x)}}
[R(x,y) - R(x,y')] > 0.
\end{align*}

Furthermore, we follow the iterative updating paradigm \citep{Deepseek-R1} by periodically resetting $\pi_\text{ref}$ to the current policy state (Algo.~\ref{alg:armor}, Line~\ref{line:ref}), avoiding saturation due to a fixed anchor.

% Empirically, combining Anchor Rollout with Mixed Optimization, which optimizes the mixture objective that takes both models into account, yields a significantly higher performance ceiling (Fig.~\ref{fig:teaser_perf}). 
% We term this synergistic framework \textbf{ARMOR}.

We designate this holistic approach---integrating the stabilization of Anchor Rollout with Mixed Optimization that calibrates the IS ratio---as \textbf{ARMOR}. 
As demonstrated in Fig.~\ref{fig:teaser_perf}, this combination yields a \textbf{significantly higher performance ceiling}.

\subsection{Explanation: Adaptive Trust Region}\label{sec:theory}
To further elucidate how Mixed Optimization enables superior asymptotic performance, we analyze its impact on policy update dynamics through the lens of the trust-region.

PPO variants (\eg GRPO, DAPO) clip the positive/negative sample's IS ratio into fixed intervals $[0, 1+\epsilon]$ or $[1-\epsilon,\infty)$ to enforce a trust region $\pi_\theta \approx \pi_{\theta_\text{old}}$. 
However, substituting our mixed ratio yields a flexible boundary.
Considering the clipping boundaries $r_{i,t}^\text{mix}(\theta) = 1 \pm \epsilon$, we have:
$$
\frac{\alpha\cdot\pi_\theta + (1-\alpha)\cdot\pi_{\text{ref}}}{\alpha\cdot\pi_{\theta_{\text{old}}} + (1-\alpha)\cdot\pi_{\text{ref}}} = 1 \pm \epsilon.
$$
Solving for the current policy $\pi_\theta$, we derive the \textbf{effective trust region boundary} \citep{TRPO,PPO}:
\begin{equation}\label{eqn:clip_expand}
\pi_\theta = \underbrace{(1 \pm \epsilon)
\cdot\pi_{\theta_{\text{old}}}}_{\text{Standard Boundary}} \pm \underbrace{\epsilon (1-\alpha)/{\alpha}\cdot\pi_{\text{ref}}}_{\text{Expansion Term }(\ge 0)}.
\end{equation}
This derivation reveals a critical mechanism: In contrast to the standard clip boundary $\pi_\theta =(1 \pm \epsilon)\cdot\pi_{\theta_{\text{old}}}$, ARMOR \textbf{dynamically expands the clipping range} proportional to the reference probability $\pi_{\text{ref}}$.
This derivation reveals that ARMOR dynamically modulates the update magnitude based on the reference prior, enabling a dual mechanism for performance breakthroughs:

\textit{1) Reinforcing verified correctness:} For positive updates ($1+\epsilon$), the higher upper bound allows larger steps towards correct actions supported by the reference. 
This is crucial for effectively learning from \emph{off-policy} correct samples, ensuring they are not unfairly clipped due to distribution shift.

\textit{2) Rectifying reference biases:} For negative updates ($1-\epsilon$), the lowered boundary permits stronger penalization of incorrect actions, even if they are highly probable in $\pi_{\text{ref}}$. 
This mechanism is key to \emph{surpassing} the reference capability, as it allows the model to decisively correct the reference's inherent biases rather than blindly imitating them.

\begin{table*}[t]
\centering
\caption{Comparison of RL baselines (DAPO, QAE) against their ARMOR continual-training counterparts across Qwen2.5-Math-7B and Qwen3-8B-Base.
ARMOR consistently improves mathematical reasoning (\texttt{avg@32}) on AIME and AMC benchmarks while maintaining comparable general capabilities (\texttt{avg@8} on GPQA and MMLU-Pro), effectively mitigating the over-optimization observed in baselines.}
\label{tab:main}
\resizebox{\linewidth}{!}{
\begin{tabular}{@{}ccc|llll|ll@{}}
\toprule
\multirow{2}{*}{\textbf{Model}} & \multirow{2}{*}{\textbf{Method}} & \multirow{2}{*}{\textbf{Step}} & \multicolumn{4}{c|}{\textbf{Math Reasoning Tasks}} & \multicolumn{2}{c}{\textbf{General Tasks}} \\ \cmidrule(l){4-9} 
 &  &  & \multicolumn{1}{c}{\textbf{AIME24}} & \multicolumn{1}{c}{\textbf{AIME25}} & \multicolumn{1}{c}{\textbf{AMC}} & \multicolumn{1}{c|}{\textbf{Average}} & \multicolumn{1}{c}{\textbf{GPQA}} & \multicolumn{1}{c}{\textbf{MMLU-Pro}} \\ \midrule
\multirow{4}{*}{\begin{tabular}[c]{@{}c@{}}Qwen2.5-\\ Math-7B\end{tabular}} & DAPO & 140 & 37.13 & 15.21 & 69.39 & 40.58 & 38.26 & 43.93 \\
 & \textbf{+ARMOR} & +200 & 43.04\up{5.91} & 18.13\up{2.92} & 76.13\up{6.74} & 45.77\up{5.19} & 42.49\up{4.23} & 46.03\up{2.1} \\ \cmidrule(l){2-9} 
 & QAE & 230 & 39.79 & 15.96 & 73.53 & 43.09 & 40.78 & 43.30 \\
 & \textbf{+ARMOR} & +240 & 41.58\up{1.79} & 16.15\up{0.19} & 76.62\up{3.09} & 44.78\up{1.69} & 39.02\down{1.76} & 42.15\down{1.15} \\ \midrule
\multirow{2}{*}{\begin{tabular}[c]{@{}c@{}}Qwen3-8B-\\ Base\end{tabular}} & DAPO & 70 & 36.98 & 28.13 & 71.72 & 45.61 & 49.87 & 65.55 \\
 & \textbf{+ARMOR} & +220 & 48.13\up{11.15} & 34.48\up{6.35} & 80.35\up{8.63} & 54.32\up{8.71} & 56.25\up{6.38} & 68.90\up{3.35} \\ \bottomrule
\end{tabular}
}
\end{table*}


\begin{paperbox}
ARMOR resolves the stability-exploration dilemma by decoupling the constraints: 
\textbf{Anchor Rollout} prevents \textbf{instability} by actively retaining correct reference solutions, while \textbf{Mixed Optimization} avoids \textbf{stagnation} by constructing an adaptive trust region. 
Empirically, ARMOR secures continuous reasoning improvements without the fragility of standard KL regularization.
\end{paperbox}

% While Anchor Rollout ensures stability, addressing the stagnation caused by uniform KL penalties requires enabling safe exploration.
% We achieve this by reformulating the optimization objective, observing that our anchor rollout process actually constructs \textbf{mixture policies} regularized by the reference policy:
% \begin{equation}
% \begin{gathered}
% \pi^{\text{mix}}_\theta = \alpha\pi_\theta + (1-\alpha)\pi_{\text{ref}},\\
% \pi^{\text{mix}}_{\theta_\text{old}} = \alpha\pi_{\theta_\text{old}} + (1-\alpha)\pi_{\text{ref}}
% \end{gathered}
% \end{equation}
% where $\alpha \in (0, 1)$ is the mixing coefficient.
% During optimization, we integrate this mixture into the loss by simply replacing $\pi_\theta$ and $\pi_{\theta_\text{old}}$ with their mixed variants.
% Consequently, the Importance Sampling (IS) term in the policy gradient framework (e.g., DAPO) becomes:
% % We integrate this mixture into the policy gradient framework (e.g., DAPO) by simply replacing the standard Importance Sampling (IS) ratio with a mixed formulation:
% \begin{equation}\label{eqn:mixed_is}
% r_{i,t}^\text{mix}(\theta) = \frac{\alpha\pi_\theta(y_{i,t}|x,y_{i,<t}) + (1-\alpha)\pi_{\text{ref}}(y_{i,t}|x,y_{i,<t})}{\alpha\pi_{\theta_{\text{old}}}(y_{i,t}|x,y_{i,<t}) + (1-\alpha)\pi_{\text{ref}}(y_{i,t}|x,y_{i,<t})}.
% \end{equation}
% Intuitively, optimizing this mixed objective can explicitly account for the anchor distribution. 
% Since $\pi_\text{ref}$ remains fixed during the update step, any performance gain in the mixture target $\pi^{\text{mix}}_\theta$ strictly translates to improvements in the ultimate target policy $\pi_\theta$:
% \begin{align*}
% &\quad\mathbb E_{y\sim\pi_{\theta}^\text{mix}(\cdot|x), y'\sim\pi_{\theta_\text{old}}^\text{mix}(\cdot|x)}
% [R(x,y) - R(x,y')] > 0 \\
% &\iff \mathbb E_{y\sim\pi_{\theta}(\cdot|x), y'\sim\pi_{\theta_\text{old}(\cdot|x)}}
% [R(x,y) - R(x,y')] > 0.
% \end{align*}
% Furthermore, we follow the iterative updating paradigm \citep{Deepseek-R1} by periodically resetting $\pi_\text{ref}$ to the current policy state (Algo.~\ref{alg:armor}, Line~\ref{line:ref}) to prevent saturation due to a fixed anchor.
% As shown in Fig.~\ref{fig:teaser_perf}, the combination of anchor rollout and mixed optimization, \ie ARMOR, \textbf{further pushes the model's performance ceiling beyond simple stabilization}, effectively utilizing signals from both the exploration policy and the anchor.

% This derivation reveals that ARMOR modulates the update magnitude based on the reference prior $\pi_{\text{ref}}$, enabling a dual mechanism to break performance bottlenecks:
% \textit{(1) Reinforcing Consensus:} For positive updates (upper bound), the expanded limit allows larger steps on correct actions supported by the reference. This amplifies the consensus on correct reasoning, accelerating convergence on verified solutions.
% \textit{(2) Surpassing the Reference:} Crucially, for negative updates (lower bound), the boundary is lowered further. This enables the model to apply stronger penalties to incorrect actions even if they are favored by the reference, allowing $\pi_\theta$ to correct the reference's inherent errors rather than blindly imitating them.
% Unlike a uniformly wider clip, this structure-aware expansion ensures the model builds upon the reference's strengths while actively unlearning its weaknesses.

% This derivation reveals that ARMOR modulates the update magnitude bi-directionally based on the reference prior $\pi_{\text{ref}}$, leading to two critical effects:
% \begin{itemize}
%     \item \textbf{Amplified Learning for Correctness (Upper Bound):} When the update is positive ($1+\epsilon$), the expansion term increases the upper limit. This allows the policy to take larger update steps towards verified correct solutions that are supported by the reference, effectively accelerating convergence.
%     \item \textbf{Aggressive Correction of Reference Errors (Lower Bound):} Conversely, when the update is negative ($1-\epsilon$), the boundary is lowered further. This enables the model to apply a stronger penalty to incorrect actions—even if they are highly probable in $\pi_{\text{ref}}$—thereby effectively bypassing the inherent mistakes or biases of the reference policy.
% \end{itemize}
% While a uniformly wider clip range could risk collapse, this adaptive mechanism ensures that large updates are structure-aware: it reinforces consensus on correct reasoning while aggressively penalizing the reference's ``blind spots''.


% this expansion is spatially selective:
% It relaxes the constraint specifically in regions where $\pi_{\text{ref}}$ is high, \ie regions grounded by the reference policy.
% This allows ARMOR to \textbf{safely incorporate larger updates} towards correct solutions while maintaining conservative elsewhere, effectively bypassing the stagnation often observed with rigid KL penalties (Sec.~\ref{sec:kl}).

% This allows the expanded boundary in ARMOR to \textbf{safely incorporate larger update ranges,} which empirically drives the model toward a higher performance ceiling, avoiding the stagnation observed with standard KL penalties (Sec.~\ref{sec:kl}).
